\documentclass[a4paper,10pt]{article}
\usepackage[utf8]{inputenc}
\usepackage[T1]{fontenc}
\usepackage[french]{babel}
\usepackage{amsmath,amssymb}
\usepackage{geometry}
\usepackage{array,tabularx}
\usepackage{tikz}
\geometry{margin=1.6cm}
\pagestyle{empty}
\newcommand{\ligne}{\par\vspace{0.55cm}\noindent\dotfill\par}
\newcommand{\carreaux}[1]{\begin{center}\begin{tikzpicture}\draw[step=0.5cm,gray!35,very thin] (0,0) grid (17,#1);\end{tikzpicture}\end{center}}
\newenvironment{exercice}[2]{\paragraph{Exercice #1 \hfill #2 points}\mbox{}\par}{\medskip}
\begin{document}
\begin{center}\Large\bfseries Devoir surveillé 10 - Etude d une série statistiques\end{center}
\vspace{2mm}
\begin{exercice}{1}{4}Définir moyenne pondérée, médiane et étendue.\end{exercice}\begin{exercice}{2}{5}Notes : $6;8;8;10;11;12;12;13;15;17$. Calculer moyenne, médiane et étendue.\end{exercice}\begin{exercice}{3}{5}Dans un tableau, les valeurs $1,2,3,4,5$ ont pour effectifs $3,7,8,5,2$. Calculer l'effectif total et la moyenne.\end{exercice}\begin{exercice}{4}{6}Comparer deux séries à partir de leurs médianes et étendues : classe A médiane $12$, étendue $9$ ; classe B médiane $11$, étendue $4$. Rédiger une interprétation.\end{exercice}\end{document}
